% Input of tikz opticla elements
\input{tikz/tikz_optical_elements.tex}

%* Da aggiungere anche qualcosa sulla fibra kagome. Magari mettere le imamgini SEM e l'immagine con le misure di core, claddiing kagome e jacket.
%* Forse da aggiungere il setup ottico nella breadboard superiore

\section{Optical dipole trap}

This section describes the implementation of the Optical Dipole Trap (ODT) setup used to transfer atoms from the MOT into a conservative optical potential. The ODT beam is guided through a Hollow-Core Photonic Crystal Fiber (HCPCF). As mentioned in Section~\ref{sec:cap_2_MOT}, the HCPCF integrated into the experimental apparatus was already present prior to this thesis work. However, a test HCPCF was employed during the development of the setup in order to optimize the alignment procedure.

During the implementation of the system, dedicated diagnostic tools were also developed to investigate the near-field beam profile and improve the fiber-coupling alignment. The following subsections describe the ODT optical configuration and the experimental procedures adopted for the alignment and loading of the trap.

%* ##########################################################
\subsection{1064nm laser characterization}
%* ##########################################################
The ODT is realized by employing the laser IPG Photonics YLR-20-1064-LP-SF, which is a Ytterbium fiber laser that emits 1064~nm linearly polarized light. The mode field diameter of the laser changes significantly upon changes of the nominal power set, therefore the mode has been characterized as a function of the power set to allow a correct realization of the HCPCF injection setup. The calibration will also provide a reference for future fine tunings of the HCPCF injection.

{\color{blue} TO BE CONTINUED, it will be moved before the optical setup}

%* ##########################################################
\subsection{Optical setup}\label{sec:optical_setup}
%* ##########################################################
In this experiment the ODT light is guided through a HCPCF, which introduces a critical complication in the injection procedure. Indeed, in contrast with standard telecom fibers, HCPCFs are multi-mode fibers. As a consequence, even small deviations of the injected beam from the ideal injection angle excites higher order modes in addition to the fundamental one. This results in distorted output beam profiles, which are not well suited for the realization of an optical conveyor belt based on the superposition of the beam guided by the HCPCF and a beam guided through a conventional telecom fiber supporting only the fundamental $LP_{01}$ mode.

For this reason, particular cure was devoted to the alignment procedure. A significant improvement in the quality of the transmitted mode was achieved by optimizing the injection while inspecting the near-field of the transmitted beam, rather than simply maximizing the transmission efficiency. The following sections describe the injection setup and the horizontal microscope built for the near-field inspection.

\subsubsection{Injection setup}
The ODT is realized by employing the laser IPG Photonics YLR-20-1064-LP-SF, which is an Yitterbium fiber laser that emits 1064~nm linearly polarized light at about 3~W of power.

\begin{figure}[h]
    \centering
    \input{tikz/odt_injection_setup.tex}
    \caption{Scheme of the optical setup for 1064~nm ODT beam  injection into the HCPCF.
    The Acousto-Optic Modulator (AOM) operates at 110~MHz, and the subsequent setup utilizes the fist diffraction order. The order zero is wasted on a beam blocker (not drawn). 
    The 780~nm light setup is not drawn, its presence will be important in the future stages of the experiment for the characterization of the atoms inside the HCPCF core.
    Fiber 2 injection setup is not drawn as well.}
    \label{fig:ODT_injection_setup}
\end{figure}

Figure \ref{fig:ODT_injection_setup} shows the optical setup scheme used for the injection of the laser light in the HCPCF.
The scheme is composed of the following elements:
\begin{itemize}
    \item \textbf{Half Wave Plate 1} (HWP1) and \textbf{Optical Isolator}. The optical isolator prevents back reflections from propagating towards the laser cavity. High-power lasers are particularly sensitive to back reflections, which can introduce power instabilities and potentially damage the laser system. Since optical isolators operate correctly only for a specific input polarization, HWP1 is used to optimize the polarization of the incoming beam and maximize the transmission through the isolator.
    
    \item \textbf{HWP2} and \textbf{Polarization Beam Splitter 1} (PBS1). These components are used to regulate the optical power delivered to the subsequent stages of the setup. The unwanted fraction of the beam is directed onto a beam block. This configuration is particularly useful during alignment procedures, where the optical power is reduced to a few milliwatts for safety reasons.
    
    \item \textbf{HWP3} and \textbf{PBS2}. PBS2 splits the beam into two different branches. The first is employed for the HCPCF injection, while the second will be injected into a PANDA polarization-maintaining fiber, which eventually will provide the second beam entering the chamber from above for the realization of the optical conveyor belt. HWP3 is adjusted to have approximately 50/50 power splitting ratio, although this value can be modified if required in future stages of the experiment.
    
    \item \textbf{Acousto-Optic Modulator} (AOM). The AOM is employed as a fast switch for the ODT beam. The device operates using an acoustic wave at \(110~\mathrm{MHz}\), causing the first diffraction order to be frequency shifted by the same amount. Although this shift is negligible with respect to the trap depth of the ODT, it will become relevant in the future implementation of the optical conveyor belt.

    \item \textbf{Beam expander}. A $\sim$3X beam expander is employed to increase the beam diameter and collimate it before the focusing conditions for coupling into the fiber.
    
    \item \textbf{Mirror 1} (M1) and \textbf{Dichroic Mirror}. Both mirrors are mounted on tilt stages and their angular adjustment is used for optimizing the fiber injection alignment. The dichroic mirror also combines the 1064~nm ODT beam with a 780~nm beam that will be used in the future for the characterization of the atoms trapped inside the HCPCF core.
    
    \item \textbf{Focusing Lens} (f50). A lens with focal length $f=50~\mathrm{mm}$ focuses the beam into the HCPCF. It is mounted on a horizontal translation stage, providing an additional degree of freedom for optimizing the injection alignment.
    
    \item \textbf{Fiber 1} (HCPCF). The HCPCF is mounted on a $x$-$y$ translation stage and terminated with an FC/APC connector. The translation stage allows transverse displacement of the fiber, providing additional alignment degrees of freedom during the injection procedure, although these degrees of freedom were found to have limited impact on the final optimization.
\end{itemize} 

{\color{blue} TO BE CONTINUED, here add the calculated mode radius at fiber injection, using results from the laser section}

\subsubsection{Horizontal microscope}

Initially, the injection into the HCPCF installed inside the vacuum chamber was attempted by using a Complementary Metal-Oxide-Semiconductor (CMOS) camera to monitoring the beam profile. However, despite the effort, strongly deformed mode profiles were obtained.

In order to better investigate the injection procedure and simplify the alignment process, a separate test HCPCF was employed outside the vacuum chamber. The test fiber consisted of an approximately 50~cm HCPCF segment with the same FC/APC connector as the fiber in the chamber. The test fiber allowed an easy measurement of the efficiency on the exit facet of the beam without any window in the path, and being outside the vacuum chamber allow for easier mode imaging. Moreover, since the two fibers shared the same connector, the optimal alignment conditions identified with the test fiber were expected to provide a good starting point for the alignment of the HCPCF installed inside the chamber. Nevertheless, even with the test fiber, optimization based solely on transmission efficiency and far-field inspection still resulted in poor beam quality and low efficiency of $\sim30\%$.

The fiber manufacturer suggested that the optimal injection condition  should be determined by inspecting the shape of the near-field of the transmitted mode. For this purpose, a horizontal microscope system was developed to inspect the HCPCF facet and to investigate the near-field profile of the transmitted mode. The optical scheme of the microscope is shown in Figure \ref{fig:horizontal_microscope}.

\begin{figure}[t]
    \centering
    \input{tikz/horizontal_microscope.tex}
    \caption{Scheme of the optical setup for the horizontal microscope.}
    \label{fig:horizontal_microscope}
\end{figure}

The fiber was positioned inside a $90^\circ$ V-groove of depth $350~\mu\mathrm{m}$. The fiber was fixed inside the groove by means of a thin foam layer attached to a tiltable mount using adhesive tape. The coating was removed from the fiber end, leaving approximately 15~mm of suspended fiber in order to allow lateral illumination of the facet.

The HCPCF facet was illuminated by a red LED focused by an $f = 50~\mathrm{mm}$ lens onto the final suspended portion of the fiber. Part of the red light propagated through the silica structure forming the kagome lattice, making them visible on the microscope camera. The observation of the kagome structure was used to optimize the microscope focus and to inspect the integrity and cleanliness of the fiber facet.

The light emerging from the fiber was collected by an optical microscope objective lens (Olympus PLN10X objective, nominal focal length $f = 18~\mathrm{mm}$, working distance $d = 10.6~\mathrm{mm}$) and subsequently re-imaged onto a CMOS camera by means of an $f = 500~\mathrm{mm}$ lens. The expected magnification is $M=f_\text{tube}/f_\text{lens}\approx27.8$. The objective is mounted on a horizontal translation stage used for the focussing. The camera was connected to a computer for real-time imaging of the transmitted near-field mode. In addition, a powermeter could be inserted between the fiber output and the objective lens in order to evaluate the coupling efficiency by comparing the transmitted power with the optical power measured before the injection stage. Finally, a black paper tube was placed between the imaging $f = 500~\mathrm{mm}$ lens and the camera to reduce background light from the LED illumination and from ambient room light. 

The alignment procedure was carried out by optimizing the transmitted near-field profile in order to obtain a mode as close as possible to a single-lobed structure, results and further details are found in Section \ref{sec:injection_alignment_nf_imaging}.

Figure~\ref{fig:photo_injection_setup_topdown} shows the top-down photographs of the complete experimental apparatus, including both the injection setup and the horizontal microscope system, and figure~\ref{fig:photo_microscope_detail_topdown} shows a detailed view of the horizontal microscope.

\begin{figure}[!h]
    \centering
    \includegraphics[angle=90, width=0.7\linewidth]{Cap_2_figures/setup_topdown.png}
    \caption{Top-down photograph of the complete experimental apparatus. The horizontal microscope, shielded by the black paper tube, is visible at the top of the image. The optical injection setup is located in the center, while the RF amplifier driving the AOM is placed at the bottom. For reference, the distance between the screw holes on the breadboard is 2.5~cm.}
    \label{fig:photo_injection_setup_topdown}
\end{figure}
\begin{figure}[!h]
    \centering
    \includegraphics[width=0.7\linewidth]{Cap_2_figures/microscope_topdown.jpg}
    \caption{Detailed top-down view of the horizontal microscope. The setup highlights the test fiber, the red LED providing lateral illumination to the fiber facet, the microscope objective lens and the refocusing lens.}
    \label{fig:photo_microscope_detail_topdown}
\end{figure}

\subsubsection{Chamber fiber optical imaging setup}
\label{sec:chamber_fiber_imaging_setup}
After the optimization and the validation of the injection alignment procedure using the test fiber, the same approach was applied to the fiber inside the vacuum chamber.

To investigate the near-field profile, a dedicated optical imaging setup was build on a breadboard positioned above the vacuum chamber (Figure~\ref{fig:upper_breadboard_imaging_setup}).
Since the optical fiber is located inside the chamber, a first stage microscope was designed to create a virtual image of the fiber facet outside the chamber.
For this purpose a f150 lens ($f=150~\mathrm{mm}$) was placed approximately 150~mm above the fiber facet, directly above the chamber window, such that the light emerging from the fiber is collimated. The collimated beam propagate through an opening and is redirected horizontally by a mirror (M1) tilted at $45^\circ$. A second lens f300 ($f=300~\mathrm{mm}$) refocuses the beam, forming an inverted image of the fiber facet with a magnification of $\times2$.

A second imaging stage, consisting of lenses L3 and L4, acted as a microscope for this intermediate image and projected a magnified image onto the sensor of a camera. The focal lengths of L3 and L4 were modified several times throughout the experiment to optimize the field of view and magnification. The configuration used most frequently employed an L3 lens with $f=100~\mathrm{mm}$ and an L4 lens with $f=500~\mathrm{mm}$, resulting in an additional magnification of $\times5$. Combined with the first imaging stage, this yielded an overall magnification of approximately $\times10$, which was used for the majority of the measurements.

\begin{figure}[!h]
    \centering
    \input{tikz/upper_breadboard_optical_setup.tex}
    \caption{Scheme of the optical setup used to image the fiber inside the vacuum chamber. The optical path between the chamber window and the mirror M1 is vertical, whereas the rest of the setup is arranged horizontally on a breadboard mounted above the chamber. The focal length of lenses L3 and L4 were varied several times during the experiment. Most images where acquired using an $f=100~\mathrm{mm}$ and an $f=500~\mathrm{mm}$ providing an overall magnification of approximately $\times 10$ combined with the first imaging stage.}
    \label{fig:upper_breadboard_imaging_setup}
\end{figure}


%* ##########################################################
\subsection{Injection alignment and near-field imaging}\label{sec:injection_alignment_nf_imaging}
%* ##########################################################
For a conventional single-mode telecom fiber, the alignment procedure is straightforward: it is sufficient to inject a low-power beam of few milliwatts, monitor the transmitted power at the fiber exit, and adjust the injected beam to maximize the guided optical power.

By contrast, aligning the injection of a multimodal HCPCFs is significantly more complex. The silica jacket that surrounds the kagome core can guide the light with coupling efficiencies as high as $60\%$, although with badly shaped mode (see figure \ref{fig:bad_injection}). Considering also that the jacket is large compared to the core, monitoring just the coupling efficiency is particularly ineffective, as the optimization can easily get trapped in a local maximum associated with jacket guiding.\\
As suggested by the manufacturer, the best procedure for the injection alignment resulted to be monitoring the near-field mode shape at the fiber exit. The complete alignment procedure begins with the use of a powermeter to ensure the impinging beam hits the fiber, even on the jacket. Then switch to live monitoring of the near-field with a camera and optimize its shape by tuning the degrees of freedom of the injection setup. Coupling efficiency is measured only after the mode profile has been optimized.

Applying this protocol, the test fiber injection was optimized  employing the injection optical setup and the horizontal microscope described in Section~\ref{sec:optical_setup}. This approach yielded to a good result and thus was subsequently applied to align the injection into the fiber in the vacuum chamber.

\begin{figure}
    \centering
    \includegraphics[width=0.5\linewidth]{Cap_2_figures/jacket_injected_fiber.png}
    \caption{Example of the near-field of a mode guided by the test fiber jacket. The light is mainly guided by the silica jacket surrounding the kagome structure, while only a weak signal is visible in the core. Note that this strongly deformed mode had a coupling efficiency of $\sim 60\%$.}
    \label{fig:bad_injection}
\end{figure}

\subsubsection{Test fiber}
The test fiber was cleaved prior to inspection.
Since the horizontal microscope has a sufficient resolution to resolve the kagome structure of the fiber, it was exploited for the focusing procedure. By monitoring the camera imaging in real-time, the kagome structure was focussed using the translation stage of the microscope objective. The objective refractive index can be considered roughly constant across the visible spectrum and into the near-infrared (NIR). This guarantees that when the structure (illuminated by red light) is in focus, the guided mode (1064~nm NIR light) is also in focus. Figure \ref{fig:kagome_images} shows an image of solely the kagome structure alongside an image of the kagome structure together with some guided light. The image show that the kagome structure is intact confirming that the cleaving was correctly performed. Furthermore, the mode is clearly guided by the core and exhibits a single lobe, however, the shape is not symmetric with respect to the fiber axis, indicating the excitation of higher order $LP_{11}$ mode alongside the fundamental $LP_{01}$. 

\begin{figure}[h]
    \centering
    \begin{subfigure}[b]{0.48\linewidth}
        \centering
        \includegraphics[width=\linewidth]{Cap_2_figures/cool_kagome.png}
    \end{subfigure}
    \hfill
    \begin{subfigure}[b]{0.48\linewidth}
        \centering
        \includegraphics[width=\linewidth]{Cap_2_figures/cool_kagome_wmode.png}
    \end{subfigure}
    \caption{Images taken with a preliminary horizontal microscope setup of the HCPCF kagome structure (left) and of the transmitted mode together with the kagome structure, with a pseudo color map (right). These images were taken with a preliminary setup of the horizontal microscope that provided a more homogeneous illumination, with the red light incident from two direction, making the kagome lattice distinctly visible.}
    \label{fig:kagome_images}
\end{figure}

The mode seen in Figure \ref{fig:kagome_images} right was associated with a transmission efficiency of $63\%$. By optimizing the alignment using the horizontal microscope the efficiency the mode shape was optimized, and an efficiency of $80\%$ has been achieved. For this optimization two screw of the mirror mounts were adjusted at the same time as a conditions walk-off, therefore performing an optimization remaining as long as possible in a good injection condition. The resulting mode is shown in Figure \ref{fig:80_eff_kagome_2} left. Figure \ref{fig:80_eff_kagome_2} right displays another mode obtained with nearly $80\%$ of injection efficiency, where the surrounding kagome structure remains visible.

\begin{figure}[h]
    \centering
    \begin{subfigure}[b]{0.48\linewidth}
        \includegraphics[width=\linewidth]{Cap_2_figures/80_mode.png}
    \end{subfigure}
    \hfill
    \begin{subfigure}[b]{0.48\linewidth}
        \includegraphics[width=\linewidth]{Cap_2_figures/10W__3.png}
    \end{subfigure}
    \caption{Images captured with the final horizontal microscope setup. A mode with $\sim 80\%$ of transmission efficiency (left) and a guided mode with the kagome structure visible (right).}
    \label{fig:80_eff_kagome_2}
\end{figure}

A final optimization was performed without the ocular lens by translating the objective slightly past the working distance and placing the camera approximately 2~m away. With this configuration, the magnification achieved is about $\sim \times 70$ (more than doubling the magnification of the horizontal microscope).
At this distance, the light scattered from the kagome structure becomes very faint, complicating the focusing procedure as the signal approaches the noise level. However, employing the FLIR camera, which has a lower read noise and a high sensitivity, the kagome lattice remained resolvable, as shown in Figure \ref{fig:kagome_hexagon_flir}.
This figure shows clearly that the near-field mode assumes a hexagonal shape, perfectly fitting the core: as expected the field inherit the symmetry of the medium in which it is confined.
The advantage of this extended projection distance is the significant expansion of the mode, the moiré patterns disturb the shape far less, allowing finer spatial details to be resolved, further improving the alignment procedure.
With this method an efficiency of $94\%$ was achieved. The resulting mode (Figure \ref{fig:mode_wo_microscope} left) is still slightly asymmetric, indicating a residual superposition of higher order modes with the fundamental one.

\begin{figure}[h]
    \centering
    \includegraphics[width=0.48\linewidth]{Cap_2_figures/flir.png}
    \caption{Images obtained with the FLIR camera using the microscope objective without the ocular lens. The mode clearly exhibits a hexagonal shape. Thanks to the FLIR high sensitivity and low read noise, the kagome structure remains visible even in this configuration. Artifacts in the top-left corner are dust grain deposited on the camera sensor.}
    \label{fig:kagome_hexagon_flir}
\end{figure}

Finally, a sweep of laser nominal power was conducted, which changes also the laser beam waist. At every tested power the injection alignment was re-optimized, and the best mode profile was obtained at 6.5~W, with an efficiency of $96\%$. The mode is shown in Figure \ref{fig:mode_wo_microscope} right, in this case the shape is far more symmetric than the previous iterations, which, together with the high transmission efficiency, suggests a significant high share of the fundamental $LP_{01}$ mode.
\begin{figure}[!h]
    \centering
    \begin{subfigure}[b]{0.48\linewidth}
        \centering
        \includegraphics[width=\linewidth]{Cap_2_figures/94_mode.png}
    \end{subfigure}
    \hfill
    \begin{subfigure}[b]{0.48\linewidth}
        \centering
        \includegraphics[width=\linewidth]{Cap_2_figures/best_shape_6.5W.png}
    \end{subfigure}
    \caption{Images obtained with the objective without the ocular lens. Mode with $94\%$ of injection efficiency (left). Optimized mode with nominal laser power set to 6.5~W which achieved $96\%$ of transmission efficiency (right).}
    \label{fig:mode_wo_microscope}
\end{figure}

{\color{blue} Qua manca qualcosa sulle due fotcamere, e anche un riferimento un pochino migliore alla HCPCF}

\subsubsection{Chamber fiber}
After validating the effectiveness of the alignment procedure using the test fiber, the same method was applied to the fiber installed inside the vacuum chamber.
Unfortunately, it was not possible to properly illuminate the fiber, and the kagome structure could not be visualized. Using the same 635~nm LED employed for the test fiber, several configurations of focusing lenses and LED positions were tested, however, no light was successfully coupled into the kagome structures. In addition, due to physical constraints imposed by the vacuum chamber geometry, it was not possible to directly illuminate the fiber facet from the outside.
One possible limitation might be the presence of a translucent rubber casing surrounding the fiber. The casing diffuses  the light of the LED preventing it to reach the kagome structure. Unlike the images obtained with the test fiber, in this case the jacket appears illuminated, suggesting that the diffused light propagated within the casing itself rather than reaching the fiber kagome structure.

Furthermore, the need to create a virtual image outside the chamber required the use of multiple lenses (see Section \ref{sec:chamber_fiber_imaging_setup}). While this approach enabled imaging of the fiber facet, it also introduced non-negligible chromatic aberrations. Figures \ref{fig:chamber_fiber_preliminary} left and \ref{fig:chamber_fiber_preliminary} right show respectively a preliminary image of the focused fiber and the same configuration with light guided through the fiber. In the latter case, the optical mode is in focus while the fiber facet appears out of focus.

The mode obtained after the injection alignment procedure is shown in Figure \ref{fig:final_mode}, which exhibited $\sim50\%$ of transmission efficiency. It should be noted that this efficiency values was measured after the light propagated through a chamber window, two lenses and two mirrors, all of which introduce transmission and reflection losses due to their non-ideal optical performance. Moreover, the fiber installed in the chamber is longer than the test fiber. Although propagation losses associated with the additional length are should be negligible, the increased length requires several bends in the fiber routing, which are known to induce leakage losses in optical fibers.

The final mode profile in Figure \ref{fig:final_mode} shows a clear asymmetry, indicating the presence of a superposition of fundamental $LP$ modes. Although, it should also be considered that, in order to enter the vacuum chamber, the fiber must be tightened in a ferrule by a swage-lock fitting. This mechanical stress can deform the fiber microstructure, altering its guiding properties and degrading the mode shape quality.

\begin{figure}
    \centering
    \begin{subfigure}[b]{0.48\linewidth}
        \centering
        \includegraphics[width=\linewidth]{Cap_2_figures/chamber_fiber.png}
    \end{subfigure}
    \hfill
    \begin{subfigure}[b]{0.48\linewidth}
        \centering
        \includegraphics[width=\linewidth]{Cap_2_figures/chamber_fiber_w_mod_.png}
    \end{subfigure}
    \caption{Preliminary images taken with a zoom about $\times10$ of the fiber facet in the vacuum chamber (left) and the fiber with a non-optimized mode (right). In contrast to the test fiber images, here, the kagome structure is not illuminated, whereas the outer jacket is. The mode exhibits a clear $LP_{11}$ profile. Note that focusing the mode bring the fiber out of focus, this makes the mode seem non-centered with respect to the fiber core.}
    \label{fig:chamber_fiber_preliminary}
\end{figure}

\begin{figure}
    \centering
    \includegraphics[width=0.5\linewidth]{Cap_2_figures/chamber_fiber_final_mode.png}
    \caption{Final mode profile obtained after the injection optimization (magnification about $\times10$). The mode shape is a bit deformed, meaning that it is a superposition of fundamental $LP$ modes.}
    \label{fig:final_mode}
\end{figure}

%
%! I think this is wrong
% In this case the objective cannot be placed at a focus length from the fiber tip because of the vacuum chamber. This introduces the necessity of a couple of lens to have a new image plane outside the chamber. Since the mode has diverged before reaching the fiber, a 2" lens is needed, this introduces a resolution limit of 
% \begin{equation*}
%     x = 1.22\,\frac{\lambda f}{D} = 1.22\,\frac{0.6~\mu\mathrm{m} \times 150~\mathrm{mm}}{50.8\mathrm{mm}} \simeq 2~\mu\mathrm{m}.
% \end{equation*}
% Since the kagome structures have a thickness of 750~nm, this prevents the whole optical system to resolve the kagome structure. 
%

%* ##########################################################
\newpage
\subsection{MOT to ODT loading sequence}
%! In realtà credo qua andrà anche uno per la CMOT
%* ##########################################################